\begin{longtable}[]{@{}>{\raggedright\arraybackslash}p{0.21\textwidth}>{\raggedright\arraybackslash}p{0.43\textwidth}>{\raggedright\arraybackslash}p{0.28\textwidth}@{}}
\caption{How conditions retained in the long tables carry into the panels and trends. Counts refer to release v1.1 and are panel rows unless stated otherwise. A trend is assigned only when consecutive editions each contribute one rule with the same unit, basis and scope; otherwise the row receives the named state. No trend field records the appearance or withdrawal of a limit.}\label{tab1PanelRules}\tabularnewline
\toprule
Condition kept in the long tables & Treatment in the panels & Trend state (panel rows) \\
\midrule
\endfirsthead
\toprule
Condition kept in the long tables & Treatment in the panels & Trend state (panel rows) \\
\midrule
\endhead
Several rules for one entity, food and edition (a multi-value cell) & Pesticide and additive cells keep every candidate value, rule and source record and leave the scalar value blank (194 pesticide and 32 additive cells). Contaminant category rows are keyed by analyte, category and basis, and no contaminant cell holds two rules. The veterinary-drug summary keeps printed animal groups; its member panel compares one animal at a time. & \path{incomparable_multiple_rules}: pesticide 110, additive 18, veterinary-drug member 1; \path{unresolved_multiple_animals}: veterinary-drug summary 69 \\
\addlinespace Printed qualifier or basis of expression & A contaminant basis such as dry weight or solid product (53 long-table rows) is part of the category-row key, so rows on different bases are never compared; in the member panel a basis that changes across editions blocks the trend. Pesticide qualifiers, including the temporary-limit asterisk, enter the scope signature that must agree across editions. & \path{incomparable_basis_unit_or_value}: contaminant member 18 (basis or unit); \path{incomparable_unresolved_review_or_source_qualifier}: pesticide 1 (source marker c) \\
\addlinespace Units that cannot be converted & Mass-per-mass units are converted before comparison. A volume-based value needs a density and keeps a blank normalised value (41 contaminant long-table rows). The additive grams-per-serving unit is not converted, and GMP is compared only with GMP. & \path{incomparable_unit_dimension_or_value}: contaminant category 3; \path{incomparable_units}: additive 47; \path{stable_gmp}: additive 256 \\
\addlinespace Scope text that changes (exclusion list, food label, use condition, animals covered) & Compared after Unicode and punctuation normalisation. Confirmed contaminant renames count as equal. Where a contaminant member heads its host category, that category's exclusions must also agree. & \path{incomparable_scope_text}: contaminant category 23 and member 22, pesticide 2,097, mycotoxin 1, veterinary-drug summary 53, additive 605 \\
\addlinespace Category without a counterpart in another edition & A row present in one edition carries no trend. A contaminant category that is no longer printed is not read as withdrawn; its foods are followed in the member panel, which releases a member only when at least two editions have a determinable host. \path{coverage_status} records when a printed name is present and does not date a limit. & \path{single_version}: contaminant 4, pesticide 906, mycotoxin 19, additive 732; \path{single_version_unmatched}: mycotoxin member 18; \path{single_2002_unmatched} and \path{single_2019_unmatched}: veterinary-drug summary 65 and 156; \path{incomparable_category_disappeared}: contaminant category 52 \\
\addlinespace Host category of a contaminant member not determinable & The edition cell and its host are left blank, and note gives the reason: the inherited rule excludes the member (8 cells), the member spans two printed rules (8), or another edition names the member in an exclusion list that differs from this rule's (1). Manually confirmed renames (5) and containment links (7) fix the host where the printed text does not. & Trend from the remaining editions; 10 rows carry a note \\
\addlinespace Quarantined fragment & Excluded from every long table and panel and kept in \path{quarantined_rows.csv}. Each of the 6 contaminant fragments duplicates a released rule of the same edition, analyte and value, so no contaminant cell is lost. The 6 additive fragments leave the 2011 or 2014 cell of 5 additive rows blank. & Trend from the editions present; 2 of the 5 additive rows carry one. A blank cell is not read as absence \\
\addlinespace Edition skipped between two editions that hold the rule & No trend spans the gap. & \path{incomparable_edition_gap}: pesticide 2, additive 1, veterinary-drug member 1 \\
\addlinespace Edition whose visible rule count differs from its declared total (GB 2763-2016 and GB 2763-2019) & Trends are kept, and \path{count_mismatch_editions} names the editions a trend relies on, so users can exclude them. The archived GB 2762-2005 and GB 2760-1996 editions were not extracted, and no panel spans them. & 5,745 of 8,602 judged pesticide rows are flagged \\
\bottomrule
\end{longtable}
